\section{Fractionally differentiated features}

\begin{qfmodel}[Fractional differencing weights: the binomial expansion of $(1-B)^d$, the recurrence, integer and fractional cases, signs and decay]\label{mod:fractional-differencing-weights}
\qffield{Statement} With the backshift operator $B^kX_t=X_{t-k}$, a real order $d$ gives the formal binomial series \citep[chap.~5, sec.~5.4, p.~77]{alpha-afml}
\[ (1-B)^d=\sum_{k=0}^\infty\binom dk(-B)^k,\qquad \tilde X_t=\sum_{k=0}^\infty\omega_kX_{t-k},\qquad \omega_k=(-1)^k\,\frac{\prod_{i=0}^{k-1}(d-i)}{k!}, \]
so $\omega=(1,-d,d(d-1)/2!,-d(d-1)(d-2)/3!,\dots)$, generated by $\omega_0=1$ and $\omega_k=-\omega_{k-1}(d-k+1)/k$ (secs.~5.4.1--5.4.2, pp.~77--78). For a positive integer $d$ the weights vanish beyond $k=d$ and ``memory beyond that point is cancelled'': $d=1$ gives $(1,-1,0,\dots)$, returns; $d=0$ gives $(1,0,\dots)$, the original series (p.~78).

\qffield{Conditions} For $0<d<1$ every weight after $\omega_0$ lies in $(-1,0)$; for $1<d<2$, $\omega_1<-1$ and $\omega_k>0$ for $k\ge2$ (p.~78, Figures~5.1--5.2). For $k>d$ the ratio $|\omega_k/\omega_{k-1}|=|(d-k+1)/k|<1$; the initial weights alternate in sign while $k<d+1$, and beyond that the weights keep one sign, negative when $\operatorname{int}[d]$ is even and positive when it is odd (sec.~5.4.3, p.~80). López de Prado concludes that the weights ``converge asymptotically to zero, as an infinite product of factors within the unit circle'' (p.~80).

Interpretation: the ratio tends to 1 (0.9986 at $k=1{,}000$ for $d=0.4$), so the product argument alone is not a proof; the decay is a power law, $|\omega_k|\approx k^{-(1+d)}/|\Gamma(-d)|$, which matches within 0.05 percent at $k=10^3$ and to four digits at $k=10^4$. The weights therefore fall slowly and the partial sums approach $\sum_k\omega_k=0$ slowly (0.017 after $10^4$ terms at $d=0.4$), which is why truncation matters (\hyperref[meth:fixed-width-fractional-differencing]{fixed-width window}).

\qffield{Practice} The recurrence equals the closed form for $d=0,0.4,1,1.5,2$; $d=0$ and $d=1$ give $(1,0,\dots)$ and $(1,-1,0,\dots)$; all weights after the first lie in $(-1,0)$ at $d=0.4$. López de Prado's motivation: integer differencing makes price series stationary ``at the expense of removing all memory'', while ``stationarity is a necessary property for inferential purposes'' and memory ``is the basis for the model's predictive power'' (sec.~5.2, p.~75); supervised learning needs stationary features, but ``stationarity does not ensure predictive power'' (p.~76). His literature review attributes the idea to Hosking (1981), whose fractionally differenced processes ``exhibit long-term persistence and antipersistence, hence enhancing the forecasting power compared to the standard ARIMA approach'' (sec.~5.3, p.~76).

\qffield{Pitfalls} Replacing $(d-i)$ by $(d-1)$ in the general weight is incorrect: at $d=0.4$ it gives $1,0.6,0.18,\dots$ instead of $1,-0.4,-0.12,\dots$; they dropped the variable from the limits and from ``$\forall k>d$''; and they turned López de Prado's ``this alternation of weight signs is necessary to make $\{\tilde X_t\}$ stationary'' (p.~80) into a sufficient condition (``makes \dots\ stationary''). For $0<d<1$ the weights do not alternate after $\omega_1$, and no sign pattern by itself makes a series stationary: that depends on the series (\hyperref[meth:fixed-width-fractional-differencing]{fixed-width window}, Conditions).

\qffield{See also} \hyperref[meth:fixed-width-fractional-differencing]{fixed-width fractional differencing}; \hyperref[def:return-measures]{return measures} (integer differencing of log prices).

\qffield{Further reading} Hosking (1981) and the fractional-integration literature López de Prado lists (p.~76).

\end{qfmodel}

\begin{qfmethod}[Fractionally differentiating a finite series: expanding window and weight loss, the fixed-width window (FFD), and choosing the minimum order $d$ with a unit-root test]\label{meth:fixed-width-fractional-differencing}
\qffield{Statement} A finite series $X_1,\dots,X_T$ cannot use infinitely many weights. With an \emph{expanding window}, $\tilde X_T$ uses $\omega_0,\dots,\omega_{T-1}$ and earlier points fewer, so they carry different amounts of memory; López de Prado drops the early points whose relative weight loss exceeds a tolerance $\tau$, such as 0.01 \citep[chap.~5, sec.~5.5.1, p.~80]{alpha-afml}. Even so the series drifts down, ``because values \dots\ are computed on an expanding window'' that keeps adding negative weights (pp.~80--82, Figure~5.3). The \emph{fixed-width window} (FFD) drops the weights once $|\omega_k|$ falls below a threshold $\tau$: with $l^*$ the last lag with $|\omega_{l^*}|\ge\tau$, $\tilde\omega_k=\omega_k$ for $k\le l^*$ and $0$ beyond, and $\tilde X_t=\sum_{k=0}^{l^*}\tilde\omega_kX_{t-k}$, the same weights for every point, which avoids the drift (sec.~5.5.2, pp.~82--83).

The \emph{minimum order} $d^*$ is the smallest $d$ for which the FFD series passes a stationarity test; López de Prado's recipe is to take the cumulative sum of a feature (so that some differencing is needed), compute FFD($d$) for $d\in[0,1]$, and keep the smallest $d$ whose ADF $p$-value is below 5 percent (secs.~5.6--5.7, pp.~84 and~88).

\qffield{Notation} Interpretation, with one consistent index (the source's is not): the point computed from the first $t$ observations misses the weights $\omega_t,\dots,\omega_{T-1}$, a relative weight loss $\lambda_t=\sum_{k=t}^{T-1}|\omega_k|/\sum_{k=0}^{T-1}|\omega_k|$, and is kept if $\lambda_t\le\tau$. FFD values exist for $t=l^*+1,\dots,T$.

\qffield{Conditions} The futures evidence as of 2018. López de Prado's evidence: E-mini S\&P~500 log prices (rolled with the ETF trick, daily) have an ADF statistic of $-0.34$ undifferenced and $-46.9$ as returns; the 95 percent critical value $-2.86$ is crossed near $d=0.35$, where the FFD series still has correlation 0.995 with the original, against 0.03 for returns (pp.~84--85, Figure~5.5); across 87 liquid futures, every log-price series passes at $d<0.6$, most at $d<0.3$, and some need no differencing (Table~5.1, pp.~85--87). He reads this as ``most studies have over-differentiated the series'' (p.~85) and the FFD series as ``a driftless blend of level plus noise \dots\ no longer Gaussian \dots\ however it is stationary'' (p.~83).

Interpretation: these are unit-root test outcomes on single samples, not demonstrations of stationarity. Snippet~5.4 runs the ADF test with one lag and a constant only, on FFD with $\tau=0.01$, a short window (11 weights at $d=0.4$); the ADF test is built against a unit root, not against fractional integration, and under fractional-integration theory a series integrated of order 1 and differenced by $d$ is integrated of order $1-d$, covariance-stationary only for $d>1/2$. The truncated weights also do not sum to zero, so a small multiple of the level remains (0.264 times $X_{t-l^*}$ at $d=0.4$, $\tau=0.01$; 0.036 at $\tau=10^{-5}$, computed here).

\qffield{Practice} At $d=0.4$ the FFD window keeps 11, 55 and 1,458 weights for $\tau=10^{-2},10^{-3},10^{-5}$; for an expanding window on 200 points, the weight loss of the point using $t$ observations falls from 0.48 at $t=1$ to 0.0135 at $t=100$, so $\tau=0.01$ keeps points from $t=118$. On 2,000 points of a simulated random walk, the one-lag ADF test on FFD($d$) with $\tau=10^{-4}$ rejects a unit root at 5 percent already at $d=0.2$ (statistic $-5.3$) and at $d=0.4$ ($-10.2$), although the theory above says such a series is not stationary for $d<1/2$: the outcome reflects the test, its lag choice and the truncation. Interpretation: choose $d$, $\tau$ and the test on training data only, and apply them unchanged out of sample; López de Prado chooses $d$ on the full E-mini history (\hyperref[pit:label-weight-leakage]{time-validity of labels and weights}).

\qffield{Pitfalls} The source contains two indexing slips: the weight-loss sum running to $T$ rather than $T-1$, with the dropped points named $\tilde X_1,\dots,\tilde X_{l^*}$ although $\lambda_l$ refers to point $T-l$; and the FFD range $t=T-l^*+1,\dots,T$, which should be $l^*+1,\dots,T$ (pp.~80 and~82). FFD output is not stationary in general; the stationarity claims rest on the minimum-$d$ procedure.

\qffield{See also} \hyperref[mod:fractional-differencing-weights]{fractional differencing weights}; \hyperref[meth:etf-trick]{ETF trick} (the rolled series López de Prado differentiates); \hyperref[pit:label-weight-leakage]{time-validity of labels and weights}.

\emph{Testing against fractional alternatives}. The ADF test asks only $d=1$ against $d=0$, so passing it at $d^*$ says nothing about the fractional order. MacKinnon and Nielsen work in Johansen's fractional VAR, $\Delta^dX_t=\Delta^{d-b}L_b\alpha\beta'X_t+\sum_{i=1}^k\Gamma_i\Delta^dL_b^iX_t+\varepsilon_t$. It nests Johansen's cointegrated VAR ($d=b=1$), the standard Dickey--Fuller test ($d=b=1$, one variable) and the fractional Dickey--Fuller test (one variable, any $d$, $b$). Its maximum-likelihood estimates of $d$, $b$ and the other parameters are asymptotically Gaussian, so they can be tested with $\chi^2$ likelihood ratios. The likelihood-ratio test of the cofractional rank, however, has a non-standard limit when $0.5<b\le d$, a functional of fractional Brownian motion that ``depend[s] on a real-valued parameter, $b$, which must be estimated''.

Its critical values must therefore be simulated for the estimated $b$, which their numerical distribution functions provide \citep[CREATES version of September 2010, abstract and secs.~1--2, eqs.~(1)--(2), pp.~1--4]{mackinnon-nielsen-2010-wp}. The model as analysed ``does not contain any deterministic terms'' (p.~2), so a constant or trend in prices is outside it. López de Prado's minimum-$d$ search is not this procedure: it tests a transformed series against a unit root, not the fractional order itself.

\qffield{Further reading} MacKinnon and Nielsen's published version (\emph{J.\ Appl.\ Econometrics}, 2014) and Johansen and Nielsen's estimation papers.

\end{qfmethod}
